\documentclass[../../../include/open-logic-section]{subfiles}

\begin{document}

\olfileid[id]{sth}{z}{story}
\olsection{Cerita yang Lebih Terperinci}
	
Dalam \olref[story][approach]{sec}, kita mengutip uraian Shoenfield tentang
proses pembentukan himpunan. Sekarang kita ingin menuliskan beberapa prinsip
lagi agar cerita ini sedikit lebih tepat. Berikut prinsip-prinsip tersebut:
\begin{enumerate}
\item[] \stageshier. Setiap himpunan dibentuk pada suatu tahap.
\item[] \stagesord. Tahap-tahap terurut: beberapa tahap berada
\emph{sebelum} tahap lainnya.\footnote{Sebenarnya, kita akan mengasumsikan
secara diam-diam bahwa tahap-tahap itu \emph{terurut baik}. Makna anggapan ini
dijelaskan dalam \olref[ordinals][]{chap}. Ini merupakan asumsi yang
substansial. Bahkan, dengan menggunakan teknik sangat cerdas yang berasal
dari \citet{Scott1974}, asumsi ini dapat \emph{dihindari}, lalu
\emph{diturunkan}. (Hal ini juga akan menjelaskan mengapa kita patut
menganggap bahwa terdapat tahap awal.) Kita tidak dapat membahasnya di sini;
untuk uraian lebih lanjut, lihat \citet{ButtonLT1}.}
\item[] \stagesacc. Untuk setiap tahap $S$, dan untuk sembarang himpunan yang
dibentuk \emph{sebelum} tahap~$S$: pada tahap~$S$ dibentuk suatu himpunan yang
anggotanya tepat semua himpunan tersebut. Tidak ada hal lain yang dibentuk
pada tahap~$S$.
\end{enumerate}
Ini merupakan prinsip-prinsip informal, tetapi kita akan dapat menggunakannya
untuk membenarkan beberapa aksioma teori himpunan Zermelo.

(Kita perlu menyampaikan satu peringatan. Walaupun kita akan menyajikan
beberapa aksioma yang sepenuhnya standar dengan nama yang sepenuhnya standar,
prinsip-prinsip bercetak miring yang baru saja kita sajikan tidak mempunyai
nama khusus dalam literatur. Kita sekadar memberi prinsip-prinsip itu julukan
yang kita harap dapat membantu.)

\end{document}
